\documentclass{amsart}
\numberwithin{equation}{section}

%%%%%%%packages
\usepackage{amssymb}
\usepackage{enumerate}
%\usepackage{showkeys}
\usepackage[bottom,first]{draftcopy}
\usepackage{xspace}
\usepackage{hyperref}

%%%%%%%%%Pagination
\hfuzz=15pt

%%%%%%%%%Theorems
\newtheorem{thm}{Theorem}[section]
\newtheorem{lem}[thm]{Lemma}
\newtheorem{cor}[thm]{Corollary} 
\newtheorem{sublem}[thm]{Sub-lemma}
\newtheorem{prop}[thm]{Proposition}
\newtheorem{conj}{Congettura}\renewcommand{\theconj}{}
\newtheorem{defin}[thm]{Definition}
\newtheorem{cond}{Condition}
\newtheorem{exmp}[thm]{Example}
\newtheorem{es}[thm]{Exercize}
\newtheorem{rem}[thm]{Comment}

%%%%%%%%%%mathcal
\newcommand\B{{\mathcal B}}
\newcommand\Co{{\mathcal C}}
\newcommand\D{{\mathcal D}}
\newcommand\F{{\mathcal F}}
\newcommand\Lp{{\mathcal L}}
\newcommand\M{{\mathcal M}}
\newcommand\Or{{\mathcal O}}
\newcommand\Pa{{\mathcal P}}
\newcommand\T{{\mathcal T}}

%%%%%%%%%%%%mathbb
\newcommand\A{{\mathbb A}}
\newcommand\C{{\mathbb C}}
\newcommand\E{{\mathbb E}}
\newcommand\N{{\mathbb N}}
\newcommand\R{{\mathbb R}}
\newcommand\Q{{\mathbb Q}}
\newcommand\To{{\mathbb T}}
\newcommand\V{{\mathbb V}}
\newcommand\Z{{\mathbb Z}}

%%%%%%%%%%%greek
\newcommand\ve{\varepsilon}
\newcommand\vf{\varphi}

%%%%%%%%%%%%%%%vaious
\newcommand\Id{\text{\bf Id}}

%%%%%   Figure-picture related definitions--beginning
\newenvironment{myfigure}{\begin{figure}[tb]\
\centering}{\end{figure}}

%\definecolor{lgray}{gray}{0.9}
%\definecolor{gray}{gray}{0.8}
%\definecolor{dgray}{gray}{0.7}
\setlength{\unitlength}{.8mm}
%%%%%   Figure related definitions--end


\pagestyle{myheadings}
\markboth{\centerline{\sc Carlangelo Liverani}\hskip-12pt}
{\centerline{\sc Mixing }\hskip-12pt}

\begin{document}
\title{\bfseries About mixing}
\author{Carlangelo Liverani}
\address{Dipartimento di Matematica\\
Universit\`{a} di Roma (Tor Vergata)\\
Via della Ricerca Scientifica, 00133 Roma, Italy\\
{\tt liverani@mat.uniroma2.it}}
\date{Rome, October 11, 2003}
\maketitle


\section{The setting}\label{sec:zero}

Consider an expanding $C^2$ map of the circle. Then one can easily
verify the Lasota-Yoke inequalities
\[
\begin{split}
&|\Lp h|_{L^1}\leq |h|_{L^1}\\
&|\Lp^nh|_{W_{1,1}}\leq \lambda^{-n}|h|_{W_{1,1}}+ D|h|_{L^1}
\end{split}
\]
where $|DT|>\lambda>1$ and $D=\sup_x\frac{|D^2_xT|}{|D_xT|^2}$.
From abstract functional analytic facts it follows that the essential
spectrum of $\Lp$ is contained in the disk of radius $\lambda$. 
\begin{lem}
The map is $T$ mixing.
\end{lem}
\begin{proof}
Let $\Pi_\theta$ be the 
projector on the eigenspace associated to the eigenvalue
$e^{i\theta}$. Consider $\Pi_\theta\Lp$, this 
is a fine rank operator and must be a diagonal operator since it is
norm bounded while Jordan blocks have norm that increases
polynomially. Let $\{h_j\}$ be a basis of the rank and consider
functionals $\ell_j$ such that $\ell_j(h_i)=\delta_{ij}$. Then $\Pi_\theta
h=\sum_j \alpha_j h_j$, thus $\alpha_j=\ell_j(h)$, hence
$\sum_jh_j\ell_j(h)=\Pi_\theta h=e^{-i\theta}\Pi_\theta\Lp
h=e^{-i\theta}\sum_jh_j\ell_j(\Lp h)$. That is 
$\ell_j(h)=e^{-i\theta}\ell_j(\Lp h)$. But the above means
\[
|\ell_j(h)|=|\ell_j(\Lp^n h)|\leq C\lambda^{-n}|h|_{W_{1,1}}+
C D|h|_{L^1}
\]
which, taking the limit for $n$ to infinity, implies $|\ell_j(h)|\leq
C D|h|_{L^1}$. Accordingly, $\ell_j$ belongs to the dual of $L^1$,
that is can be identified with a function $\bar \ell_j\in L^\infty$
via the formula $\ell_j(h)=\int\bar \ell_j(x)h(x) dx$ and $\bar
\ell_j\circ T=e^{i\theta}\ell_j$ .  

This means that $\Lp^n\bar \ell_j^*=e^{i\theta n}\Lp^n(\bar \ell_j^*\circ
T^n)=e^{i\theta n}\bar \ell_j^*\Lp^n 1$. Or, better,
\[
\frac 1n \sum_{k=0}^{n-1}e^{-ik\theta}\Lp^k\bar \ell_j^*=\bar
\ell_j^*\frac 1n \sum_{k=0}^{n-1} \Lp^k 1.
\]
By compactness the last averages have a subsequence converging to
$\Pi_\theta\ell_j^*$ and $\Pi_1 1$, respectively. Hence 
\[
\Pi_\theta\ell_j^*=\ell_j^*\Pi_1 1.
\]
In addition, it is easy to verify that if the density 
is zero at same point, then it must be zero on all its
preimages. Since such preimages are dense, then the density would be
zero which is absurd since it must have integral one. So the density
is strictly positive and thus, since it is continuous, uniformly
so. But his implies that $\bar\ell_j$ has a version in $W_{1,1}$.
 
Now that means that all the
sets $A_\alpha:=\{|\bar \ell_j|\leq \alpha\}$ of positive measure must
contain an interval 
$I$. But then $A_\alpha=T^n A_\alpha$ mod zero Lebesgue measure sets,
but for $n$ large enough $T^nI=\To^1$, thus $A=\To^1$. This means that
$|\bar \ell_j|$ must be constant. We
can write $\bar \ell_j=Ce^{i\vf}$ for 
some real function $\vf$ everywhere continuous but at a point $x_*$
where it has a jump that is an integer multiple of $2\pi$ and the
point $x_*$ can be chosen arbitrarily. Hence
\[
e^{i\vf\circ T}=\bar \ell_j\circ T=e^{i\theta}\bar
\ell_j=e^{i(\vf+\theta)}.
\]
that is $\vf\circ T=\vf+\theta\mod 2\pi$. Finally, notice that the map $T$ has
at least one fixed point $\bar x$, thus, choosing $x_*\neq \bar x$,
$\theta=2\pi \Theta(\bar x)$ which means that $e^{i\theta}=1$.
So non other eigenvalues can be present beside the eigenvalue
one. Moreover, since $\ell_j=\ell_j\circ T$ implies $\ell_j$ constant,
one must be a simple eigenvalue. This implies that, for all $h\in
W_{1,1}$, $\int h=1$, and $\vf\in L^\infty$, holds
\[
\lim_{n\to\infty}\int_{\To^1}\vf\circ T^n
h=\int \vf h_*+\Or(|\vf|_{L^\infty}\nu^{-n}|h|_{W_{1,1}}),
\]
where $h_*$ is the eigenvector associated to one and $\nu$ is the
radius of the spectrum of $\Lp$ once  one is 
removed. By a standard approximation argument it follows that the
system is mixing.
\end{proof}
\end{document}
